\documentclass[11pt,border=4pt]{standalone}
% Estilo comum dos esquemáticos do relatório (circuitikz).
% Uso: \documentclass[11pt,border=4pt]{standalone} e, em seguida, % Estilo comum dos esquemáticos do relatório (circuitikz).
% Uso: \documentclass[11pt,border=4pt]{standalone} e, em seguida, % Estilo comum dos esquemáticos do relatório (circuitikz).
% Uso: \documentclass[11pt,border=4pt]{standalone} e, em seguida, \input{estilo_circuitos}.
\usepackage[T1]{fontenc}
\usepackage[utf8]{inputenc}
\usepackage[brazilian]{babel}
\usepackage{lmodern}
\usepackage{amsmath}
\usepackage[american]{circuitikz}
\usetikzlibrary{calc,arrows.meta}

% Cores das figuras originais
\definecolor{anot}{RGB}{0,135,200}      % anotações: valores calculados, setas, resistências vistas
\definecolor{sombra}{RGB}{220,242,253}  % região do transistor nos modelos T e pi
\definecolor{term}{RGB}{62,124,148}     % nomes de terminais destacados (B, C, E)

% Fios com 0,4 pt e corpo dos componentes com 0,8 pt, como nas figuras originais
\ctikzset{bipoles/thickness=2}
.
\usepackage[T1]{fontenc}
\usepackage[utf8]{inputenc}
\usepackage[brazilian]{babel}
\usepackage{lmodern}
\usepackage{amsmath}
\usepackage[american]{circuitikz}
\usetikzlibrary{calc,arrows.meta}

% Cores das figuras originais
\definecolor{anot}{RGB}{0,135,200}      % anotações: valores calculados, setas, resistências vistas
\definecolor{sombra}{RGB}{220,242,253}  % região do transistor nos modelos T e pi
\definecolor{term}{RGB}{62,124,148}     % nomes de terminais destacados (B, C, E)

% Fios com 0,4 pt e corpo dos componentes com 0,8 pt, como nas figuras originais
\ctikzset{bipoles/thickness=2}
.
\usepackage[T1]{fontenc}
\usepackage[utf8]{inputenc}
\usepackage[brazilian]{babel}
\usepackage{lmodern}
\usepackage{amsmath}
\usepackage[american]{circuitikz}
\usetikzlibrary{calc,arrows.meta}

% Cores das figuras originais
\definecolor{anot}{RGB}{0,135,200}      % anotações: valores calculados, setas, resistências vistas
\definecolor{sombra}{RGB}{220,242,253}  % região do transistor nos modelos T e pi
\definecolor{term}{RGB}{62,124,148}     % nomes de terminais destacados (B, C, E)

% Fios com 0,4 pt e corpo dos componentes com 0,8 pt, como nas figuras originais
\ctikzset{bipoles/thickness=2}

% Circuito equivalente c.c. da Questão 1, com os capacitores abertos (fig:q1-cc).
% Compilar com pdflatex; o PDF vai para figuras/fig_q1_cc.pdf e entra no relatório em escala 1.
% Os rótulos são nós explícitos (inner sep=0pt) com afastamento fixo de 6 a 7 pt do elemento,
% para manter a folga mínima de 3 pt medida pelo folga.py.
\begin{document}
\begin{circuitikz}
  % ================= M1 (à esquerda): porta em 0 V por R_G =================
  \node[nmos, arrowmos] (M1) at (0,0) {};
  \coordinate (A) at ($(M1.D |- M1.G)+(-2.32,0)$);   % canto porta / R_G
  % dreno direto em +10 V
  \draw (M1.D) -- ++(0,0.47) node[vcc] (P1) {};
  % porta: R_G até a terra
  \draw (M1.G) -- (A) to[R] ++(0,-3.04) node[ground] {};
  % fonte: I_D1, nó V_S1, R_S1 e -10 V
  \draw (M1.S) -- ++(0,-1.06) coordinate (S1) to[R] ++(0,-2.80) node[vee] (N1) {};
  \node[currarrow, rotate=-90] (I1) at ($(M1.S)!0.5!(S1)$) {};
  \node[circ] at (S1) {};

  % ================= M2 (à direita): porta aterrada =================
  \node[nmos, arrowmos] (M2) at (7.60,0) {};
  % dreno: nó V_D2, R_D e +10 V
  \draw (M2.D) -- ++(0,0.71) coordinate (D2) to[R] ++(0,2.34) node[vcc] (P2) {};
  \node[circ] at (D2) {};
  % porta direto na terra
  \draw (M2.G) -- ($(M2.D |- M2.G)+(-2.08,0)$) node[ground] {};
  % fonte: I_D2, nó V_S2, R_S2 e -10 V
  \draw (M2.S) -- ++(0,-1.06) coordinate (S2) to[R] ++(0,-2.80) node[vee] (N2) {};
  \node[currarrow, rotate=-90] (I2) at ($(M2.S)!0.5!(S2)$) {};
  \node[circ] at (S2) {};

  % ================= rótulos =================
  % meia largura do zigue-zague do resistor: 0,21 cm; haste de vcc/vee: 0,42 cm
  \begin{scope}[every node/.style={inner sep=0pt}]
    % M1
    \node[anchor=south] at ($(P1)+(0,0.42cm+7pt)$) {$+10\,\mathrm{V}$};
    \node[anchor=west] at ($(M1.D)!0.5!(M1.S)+(6pt,0)$) {$M_1$};
    \node[anchor=south] at ($(A)+(0,6pt)$) {$V_{G1} = 0$};
    \node[anchor=east, align=left] at ($(A)+(-0.21cm-6pt,-1.52cm)$) {$R_G$\\$400\,\mathrm{k\Omega}$};
    \node[anchor=east] at ($(I1)+(-7pt,0)$) {$I_{D1}$};
    \node[anchor=west] at ($(S1)+(6pt,0)$) {$V_{S1}$};
    \node[anchor=east, align=left] at ($(S1)+(-0.21cm-6pt,-1.40cm)$) {$R_{S1}$\\$10\,\mathrm{k\Omega}$};
    \node[anchor=north] at ($(N1)+(0,-0.42cm-6pt)$) {$-10\,\mathrm{V}$};
    % M2
    \node[anchor=south] at ($(P2)+(0,0.42cm+7pt)$) {$+10\,\mathrm{V}$};
    \node[anchor=west, align=left] at ($(D2)+(0.21cm+6pt,1.17cm)$) {$R_D$\\$5\,\mathrm{k\Omega}$};
    \node[anchor=west] at ($(D2)+(6pt,0)$) {$V_{D2}$};
    \node[anchor=west] at ($(M2.D)!0.5!(M2.S)+(6pt,0)$) {$M_2$};
    \node[anchor=east] at ($(I2)+(-7pt,0)$) {$I_{D2}$};
    \node[anchor=west] at ($(S2)+(6pt,0)$) {$V_{S2}$};
    \node[anchor=east, align=left] at ($(S2)+(-0.21cm-6pt,-1.40cm)$) {$R_{S2}$\\$10\,\mathrm{k\Omega}$};
    \node[anchor=north] at ($(N2)+(0,-0.42cm-6pt)$) {$-10\,\mathrm{V}$};
  \end{scope}
\end{circuitikz}
\end{document}
